\chapter{Building a process monitor}
\label{ch:procmon}

\begin{aside}
htop, but built in \maya{}. Periodic sampling, sorting,
filtering, killing processes. A practical second-tier
example.
\end{aside}

\section{Sampling}

Every second, parse \code{/proc} on Linux (or \code{ps} on
others). Build a list of process records.

\begin{lstlisting}
struct Process {
    int pid;
    std::string name;
    float cpu_pct;
    int rss_kb;
    int threads;
};
\end{lstlisting}

\section{Display}

A table widget with sortable columns. PID, name, CPU,
memory, threads.

\section{Sort}

Click a header to sort by that column. Toggle direction
on second click.

\section{Filter}

A search box at the top filters by name substring.

\section{Kill}

Selected process + \code{K} key sends SIGTERM. A confirm
dialog before sending SIGKILL.

\section{Refresh rate}

User-configurable: 0.5 s, 1 s, 2 s, 5 s. Faster gives
fresher data but uses more CPU itself.

\section{Sparkline}

Each row could include a tiny CPU sparkline showing the
last 60 seconds. Use the spark widget from
Chapter~\ref{ch:charts}.

\section{Hierarchical view}

A tree view shows parent-child relationships. Group by
process group.

\section{Recap}

\begin{recap}
\begin{itemize}[leftmargin=*]
  \item Periodic sample of /proc.
  \item Sortable, filterable table.
  \item Kill with confirmation.
  \item Configurable refresh rate.
  \item Optional sparklines per row.
\end{itemize}
\end{recap}

\section*{Workshop}

\begin{enumerate}[leftmargin=*]
  \item Build the table view with five columns.
  \item Add filter and sort.
  \item Wire kill with confirm dialog.
\end{enumerate}
